\chapter{Interest Rates} \label{chap:interest_rates}

As established in the discussion of the Cost of Capital (Chapter \ref{chap:discount_rate}), particularly concerning the Risk-Free Rate (\(R_F\)) in Section \ref{sec:risk_free_rate}, interest rates form the bedrock upon which much of the Discounted Cash Flow valuation is built. The risk-free rate serves as the fundamental benchmark return in the financial market, influencing the cost of both equity and debt, and potentially guiding assumptions about long-term growth.

This chapter delves into the impact of interest rate fluctuations on the valuation outcomes derived from the DCF model framework developed in the preceding chapters. We will explore how changes in the benchmark risk-free rate, both in absolute terms and through shifts in the yield curve structure, propagate through the model's inputs to affect the final estimated intrinsic value. Furthermore, we will examine the distinct but related impact of changes in default spreads, which influence both the cost of debt for individual firms and the risk assessment for investments in specific countries. Understanding these dynamics is crucial for interpreting valuation results in the context of evolving macroeconomic conditions and market sentiment.

\section{Central Banks} \label{sec:central_banks}
Central Banks are pivotal institutions within modern economies, typically tasked with maintaining price stability (often defined as low and stable inflation, usually targeting around 2-3\%) and, in some jurisdictions like the United States, promoting maximum sustainable employment. They achieve these primary objectives by influencing the general level of interest rates and the availability of credit in the economy through various policy tools. The key policy rates set by major central banks serve as benchmarks for the entire financial system.

\paragraph{Federal Reserve (FED) - United States}
The FED sets monetary policy for the US Dollar area, primarily through the Federal Open Market Committee (FOMC). Its key rates include:
\begin{itemize}
    \item \textbf{Federal Funds Rate Target Range}: This is the target range set by the FOMC for the rate at which commercial banks lend reserves to each other overnight. Market forces determine the effective federal funds rate within this target range.
    \item \textbf{Interest on Reserve Balances (IORB)}: The interest rate paid by the Fed on reserves held by eligible institutions. This rate acts as a key administered rate, helping to keep the effective federal funds rate within the target range by providing banks with a risk-free overnight investment option at the Fed.
    \item \textbf{Overnight Reverse Repurchase Agreement (ON RRP) Facility Offering Rate}: The rate offered to a broad set of counterparties (including money market funds) for placing funds with the Fed overnight, secured by Treasury securities. This rate typically acts as a floor for short-term money market rates, reinforcing the lower bound of the target range.
    \item \textbf{Discount Rate}: The rate at which commercial banks can borrow directly from the Fed through the "discount window." It is typically set above the top of the federal funds target range (acting as a ceiling) to encourage banks to borrow from each other first, positioning the Fed as a lender of last resort.
\end{itemize}

\paragraph{European Central Bank (ECB) - Eurozone}
The ECB sets monetary policy for the Euro area. Its key interest rates are:
\begin{itemize}
    \item \textbf{Main Refinancing Operations (MRO) Rate}: The interest rate for the regular weekly liquidity-providing operations conducted by the ECB with commercial banks. This rate influences the cost for banks to borrow reserves from the ECB.
    \item \textbf{Deposit Facility Rate}: The interest rate banks receive for depositing funds overnight with the ECB. It normally acts as a floor for overnight market interest rates in the Eurozone.
    \item \textbf{Marginal Lending Facility Rate}: The interest rate at which banks can borrow overnight liquidity directly from the ECB against eligible collateral. It normally provides a ceiling for the overnight market interest rate.
\end{itemize}

\subsection{Monetary Policy Implementation}
Central Banks employ several instruments to effectively implement their monetary policy decisions and influence broader financial conditions:
\begin{itemize}
    \item \textbf{Open Market Operations (OMO)}: Historically the primary tool, OMOs involve the central bank buying or selling securities (usually government bonds) in the open market. Purchases inject liquidity (reserves) into the banking system, putting downward pressure on short-term interest rates, while sales withdraw liquidity, pushing rates up. In modern frameworks with ample reserves, OMOs are often used more structurally to maintain the desired level of bank reserves rather than for daily rate adjustments.
    \item \textbf{Administered Rates (Fed's Modern Framework)}: With banking systems often holding abundant reserves (partly due to past QE), central banks like the Fed increasingly rely on setting administered rates – the IORB and ON RRP rates – to directly influence market rates by setting the reservation rates for holding or borrowing cash. These rates effectively guide the federal funds rate within its target range.
   \item \textbf{Reserve Requirements}: The minimum percentage of deposits that commercial banks must hold as reserves (either in their vaults or at the central bank) and cannot lend out. Lowering this requirement frees up funds for banks to lend, potentially stimulating the economy, while raising it restricts lending. This tool is used less frequently today in many major economies compared to OMOs or administered rates.
   \item \textbf{Quantitative Easing (QE) and Quantitative Tightening (QT)}: QE involves large-scale purchases of assets (typically longer-term government bonds or mortgage-backed securities) by the central bank, aimed at lowering long-term interest rates, increasing liquidity, and encouraging risk-taking when short-term policy rates are at or near zero. QT is the opposite process, where the central bank shrinks its balance sheet, usually by allowing assets to mature without reinvesting the principal or by actively selling assets. QT tends to withdraw liquidity and put upward pressure on longer-term yields.
   \item \textbf{Forward Guidance and Communication}: Central banks actively use communication – through policy statements, press conferences, economic projections, and speeches – to signal their intentions regarding the future path of policy rates and other tools. By shaping market expectations, forward guidance can influence longer-term interest rates and financial conditions even without immediate changes to policy settings.
\end{itemize}
The specific mix and emphasis on these tools evolve based on economic conditions and the central bank's operational framework.

\section{Impact of Risk-Free Rate Changes on Valuation}

Changes in the prevailing risk-free rate (\(R_F\)) directly and indirectly influence the estimated value of a firm through several key components of the DCF model.

\paragraph{Direct Impact via Discount Rate (WACC)}
The most immediate effect of a change in \(R_F\) is on the Weighted Average Cost of Capital (WACC), the discount rate used to bring future cash flows back to their present value (Equation \ref{eq:dcf_formula}).
\begin{itemize}
    \item \textbf{Cost of Equity (\(C_E\))}: According to the Capital Asset Pricing Model (CAPM, Equation \ref{eq:cost_of_equity}), \(C_E = R_F + \beta_L \times ERP\). Thus, an increase in \(R_F\) directly increases the required return on equity, assuming Beta and the Equity Risk Premium remain constant.
    \item \textbf{Cost of Debt (\(C_D\))}: The pre-tax cost of debt is typically conceptualized as \(C_D = R_F + \text{Company Default Spread}\) (Section \ref{sec:cost_of_debt}). An increase in \(R_F\) will generally lead to an increase in the market rates demanded for new corporate debt, thus raising \(C_D\).
\end{itemize}
Since both the cost of equity and the cost of debt are positively related to \(R_F\), an increase in the risk-free rate will lead to a higher WACC (Equation \ref{eq:wacc}). A higher WACC, in turn, reduces the present value of both the projected Free Cash Flows (\(FCFF_n\)) during the forecast period and the Terminal Value (\(TV\)), leading to a lower estimated intrinsic value for the firm, all else being equal. Conversely, a decrease in \(R_F\) tends to lower WACC and increase the estimated firm value.

\paragraph{Indirect Impact via Perpetual Growth Rate (\(g\))}
As discussed in Section \ref{sec:estimating_g}, the long-term risk-free rate (\(R_F\)) is often used as a proxy or ceiling for the perpetual growth rate (\(g\)) used in the Stable Growth Terminal Value calculation (Equation \ref{eq:terminal_value_stable_growth}). If \(g\) is explicitly linked to \(R_F\), then a change in the risk-free rate could also imply a change in the assumed long-term growth potential.

An increase in \(R_F\) might suggest higher expected nominal GDP growth (inflation plus real growth), potentially justifying a higher \(g\). Looking at the terminal value formula, \(TV = \frac{FCFF_N \times (1+g)}{\text{WACC}_{\text{Terminal}} - g}\), an increase in \(g\) has a positive effect on the numerator and reduces the denominator spread (\(\text{WACC}_{\text{Terminal}} - g\)), both pushing the \(TV\) higher.

However, the simultaneous increase in \(\text{WACC}_{\text{Terminal}}\) (due to the higher \(R_F\)) increases the denominator spread, pushing the \(TV\) lower. The net effect on the terminal value depends on the sensitivity of the valuation to changes in the discount rate versus changes in the growth rate. Generally, the impact through the discount rate (\(\text{WACC}_{\text{Terminal}}\)) tends to be more dominant, meaning that the overall effect of a higher \(R_F\) is typically a lower terminal value, despite the potential increase in \(g\).

\paragraph{Connection to Stock Prices}
The DCF model provides a framework for estimating intrinsic value. Market prices reflect collective investor expectations and required rates of return. Therefore, sustained increases in benchmark interest rates generally lead to higher required returns (higher WACC) across the market, reducing the present value investors place on future earnings and cash flows. This translates into downward pressure on stock valuations and market prices, assuming other factors like earnings expectations remain unchanged. Conversely, falling interest rates tend to support higher stock valuations.

\subsection{Absolute vs. Relative Changes (Yield Curve)}

The impact of interest rate changes can differ depending on whether the entire yield curve shifts in parallel (absolute change) or whether its shape changes (relative change).

\paragraph{Absolute Changes (Parallel Shifts)}
A parallel shift occurs when interest rates across all maturities increase or decrease by roughly the same amount. For example, if \(R_F\) (proxied by the long-term government bond yield) increases, and short-term rates also increase similarly, this absolute rise impacts all components sensitive to \(R_F\). The Cost of Equity increases via CAPM, the base rate for the Cost of Debt increases, and thus the WACC rises across the board. If \(g\) is linked to \(R_F\), it may also adjust upwards. This scenario aligns with the general discussion above, typically leading to lower valuations.

\paragraph{Relative Changes (Yield Curve Shifts)}
The yield curve depicts interest rates across different maturities. Its shape can change:
\begin{itemize}
    \item \textbf{Steepening}: Long-term rates rise more than short-term rates (or fall less). This might reflect expectations of higher future inflation, stronger future economic growth, or increased uncertainty demanding a higher term premium.
    \item \textbf{Flattening/Inversion}: Short-term rates rise more than long-term rates (or fall less). Flattening often occurs late in economic cycles as central banks raise short-term rates, while inversion (short-term rates higher than long-term) is sometimes seen as a predictor of recession.
\end{itemize}
While the DCF model typically relies on a long-term \(R_F\) (matching the duration of the cash flows), significant changes in the yield curve's shape can still be relevant. A steepening curve might reinforce the use of a higher long-term \(R_F\) and potentially signal higher growth expectations (\(g\)). Conversely, a flattening or inverted curve might signal economic slowdown, potentially tempering long-term growth assumptions (\(g\)) even if the long-term \(R_F\) itself hasn't fallen dramatically. While the model primarily uses the selected long-term \(R_F\), the overall yield curve shape provides crucial context about market expectations that can inform the qualitative judgments underlying forecasts for growth and risk.

\section{Industries Impact of Rates Changes} \label{sec:industry_impact_rates}

While the previous section outlined the general mechanisms through which changes in risk-free rates affect firm valuations via the WACC and potential growth assumptions – a dynamic applicable across the economy – the magnitude and nuances of these effects are not uniform across all sectors and industries. Different industries exhibit varying degrees of sensitivity to interest rate fluctuations due to fundamental differences in their operating and financial characteristics.

Factors contributing to these differential impacts include:
\begin{itemize}
    \item \textbf{Capital Intensity and Investment Cycles}: Industries requiring significant, long-term capital investment (e.g., utilities, manufacturing, real estate, infrastructure) are often more sensitive to changes in long-term rates, as these directly impact the cost of financing large projects and the hurdle rates for new investments.
    \item \textbf{Leverage Levels}: Companies and industries that typically operate with higher levels of debt (e.g., financials, utilities, private equity-backed firms) will experience a more pronounced impact on their WACC and interest expense from rate changes compared to firms with low debt levels.
    \item \textbf{Demand Sensitivity}: Demand for goods and services in some sectors is highly sensitive to interest rates (e.g., housing, automobiles, durable consumer goods), as borrowing costs influence consumer purchasing power and decisions. Other sectors (e.g., consumer staples, healthcare) tend to have more inelastic demand.
    \item \textbf{Growth Profiles and Valuation Duration}: High-growth companies, particularly those whose value is heavily weighted towards cash flows expected far in the future (e.g., technology, biotechnology), often exhibit higher "duration" in their valuations. This means their present values are mathematically more sensitive to changes in the discount rate (WACC) compared to mature, stable companies with more front-loaded cash flows.
    \item \textbf{Business Model Specifics}: Certain business models are intrinsically linked to interest rates. For example, banks' net interest margins are directly affected by the spread between lending and deposit rates, which are influenced by policy rates and the yield curve shape. Insurance companies' profitability depends on investment returns earned on their float, which is sensitive to prevailing yields.
\end{itemize}

Therefore, while the general principle holds – rising rates tend to pressure valuations downwards and vice versa – a more granular, industry-specific analysis is crucial for accurately assessing the impact of monetary policy shifts. Following the micro-level approach emphasized in this thesis, the subsequent analysis would ideally consider these sector-specific sensitivities when interpreting valuation results or forecasting performance under different interest rate scenarios, keeping in mind the base framework outlined previously.

\subsection{Energy} \label{subsec:industry_energy}
The energy sector's sensitivity to interest rates can be mixed. Commodity prices (oil, gas) are often driven more by supply/demand dynamics and geopolitical events than directly by interest rates. However, energy projects (exploration, production, renewables) are highly capital-intensive, making financing costs sensitive to rate changes. Higher rates increase the cost of funding large projects and can make marginal projects less viable, potentially impacting long-term supply investment. Additionally, as a cyclical sector, energy demand can be indirectly affected by the broader economic slowdown often associated with rising rates.

\subsection{Utilities} \label{subsec:industry_utilities}
Utilities are traditionally considered one of the most interest-rate-sensitive sectors. They are highly capital-intensive (requiring significant investment in infrastructure like power plants and grids) and typically operate with substantial debt loads. Higher interest rates directly increase their financing costs and WACC. Furthermore, utility stocks are often viewed by investors as bond proxies due to their stable cash flows and relatively high dividend yields. When interest rates rise, the yields on actual bonds become more attractive, potentially leading investors to shift away from utility stocks, putting downward pressure on their valuations.

\subsection{Information Technology} \label{subsec:industry_it}
The Information Technology sector, particularly high-growth software and semiconductor companies, tends to be negatively sensitive to rising interest rates. This sensitivity stems primarily from valuation dynamics. Much of the value of these companies lies in expected cash flows far in the future (high valuation duration). Higher interest rates increase the discount rate (WACC), which disproportionately reduces the present value of these distant cash flows. While many large tech firms have strong balance sheets with low net debt, the valuation effect often dominates.

\subsection{Consumer Staples} \label{subsec:industry_consumer_staples}
Consumer Staples (food, beverages, household products) are generally considered defensive and less sensitive to interest rate changes compared to cyclical sectors. Demand for essential goods tends to be stable regardless of economic conditions or borrowing costs. While companies still use debt, their leverage is often moderate. However, like utilities, some mature staples companies with stable dividends might be viewed partially as bond proxies, experiencing some valuation pressure when rates rise significantly. Their stable cash flows, however, provide resilience.

\subsection{Health Care} \label{subsec:industry_health_care}
Health Care exhibits mixed sensitivity. Demand for healthcare services and pharmaceuticals is relatively inelastic, making the sector defensive during economic downturns often triggered by rate hikes. However, biotechnology and certain medical device companies, particularly those in early stages or reliant on future product launches, share characteristics with high-growth tech stocks (high valuation duration) and can be sensitive to discount rate changes. Established pharmaceutical and managed care companies with strong cash flows are generally more resilient.

\subsection{Materials} \label{subsec:industry_materials}
The Materials sector (chemicals, metals, mining, construction materials) is typically cyclical and somewhat sensitive to interest rates. Demand is closely tied to industrial production and construction activity, both of which can be dampened by higher borrowing costs and economic slowdowns. Companies in this sector can also be capital-intensive, making investment decisions sensitive to financing costs. Commodity price fluctuations, however, often play a more dominant role than interest rates directly.

\subsection{Industrials} \label{subsec:industry_industrials}
Industrials (machinery, aerospace, defense, transportation, construction) represent a broad, cyclical category. Sensitivity to interest rates varies. Companies involved in large capital projects (construction, heavy machinery) or selling big-ticket items sensitive to financing (aerospace) can be negatively impacted by rising rates. Transportation can be affected by economic activity levels influenced by rates. Defense spending may be less sensitive to economic cycles driven by rates. Higher rates increase financing costs for expansion and M\&A, common activities in this sector.

\subsection{Real Estate} \label{subsec:industry_real_estate}
The Real Estate sector (particularly REITs - Real Estate Investment Trusts) is highly sensitive to interest rates. Property values and development activity are directly influenced by mortgage rates and financing costs. Higher rates make borrowing more expensive for developers and buyers, potentially dampening demand and transaction volumes. For REITs, higher rates increase their cost of debt, impacting profitability. Similar to utilities, REITs are often valued for their income (dividends from rental streams), making them susceptible to competition from higher bond yields when interest rates rise.

\subsection{Consumer Discretionary} \label{subsec:industry_consumer_discretionary}
Consumer Discretionary (automobiles, retail, travel, luxury goods) is generally quite sensitive to interest rate changes. Purchases of big-ticket items like cars and expensive goods are often financed, making demand directly responsive to borrowing costs. Higher rates can also signal or contribute to economic slowdowns, reducing consumer confidence and spending on non-essential items. Retailers may also face higher inventory financing costs.

\section{Impact of Default Spread Changes}

Beyond the risk-free rate, changes in default spreads – the premium demanded by investors for bearing credit risk – also significantly impact DCF valuations.

\paragraph{Impact on Cost of Debt (\(C_D\))}
The company-specific default spread is a direct input into the Cost of Debt (\(C_D = R_F + \text{Company Default Spread}\), Section \ref{sec:cost_of_debt}). This spread reflects the market's assessment of the company's creditworthiness, often inferred from credit ratings (actual or synthetic, see Table \ref{tab:interest_coverage_ratings}) or market data like bond yields or CDS spreads. An increase in the company's perceived default risk, or a general increase in market risk aversion leading to wider spreads across the board, will increase its specific default spread. As a proxy for such general market trends and risk aversion concerning corporate credit risk, market participants often monitor credit default swap (CDS) indices. Prominent examples include the S\&P Global Markit iTraxx indices in Europe (with iTraxx Europe Main covering investment-grade and iTraxx Crossover covering sub-investment-grade entities) and the S\&P Global Markit CDX indices in North America (with similar investment-grade and high-yield benchmarks). A widening (increase) in these index spreads typically signals heightened perceived risk across the relevant corporate sector (investment grade or high yield), suggesting potential upward pressure on individual company default spreads within that category, while narrowing spreads indicate improving credit conditions. This overall effect raises the \(C_D\) for affected companies, subsequently increasing the WACC and lowering the estimated firm value. For example, if a company's financial performance deteriorates, leading to a downgrade in its credit rating, the associated higher default spread will increase its WACC; similarly, a general market panic reflected in widening iTraxx/CDX spreads would also likely increase the required \(C_D\) even if the company's specific situation hasn't changed dramatically.

\paragraph{Impact on Risk-Free Rate and ERP for Non-AAA Countries}
Changes in sovereign default spreads have a dual impact, particularly for companies operating in or exposed to countries not considered risk-free (i.e., non-AAA rated).
\begin{enumerate}
    \item \textbf{Adjusted Risk-Free Rate (\(R_F\))}: As discussed in Section \ref{sec:adjust_debt_default_risk}, when using the subtraction method (\(R_F = Y_{Govt} - \text{Country Default Spread}\)), the risk-free rate is derived by adjusting the local government bond yield. If the sovereign default spread widens (e.g., due to increased perceived country risk reflected in higher CDS spreads like those in Table \ref{tab:cds_country_prices_example} or rating downgrades leading to higher spreads like in Table \ref{tab:rating_default_spreads_example}), the resulting impact on the calculated \(R_F\) depends on the relative movements of the components. Specifically, the direction of change in \(R_F\) is determined by whether the increase (or decrease) in the local government bond yield (\(Y_{Govt}\)) is greater or smaller than the simultaneous widening of the Country Default Spread. Consequently, the adjusted \(R_F\) could either increase or decrease depending on these relative market dynamics.
    \item \textbf{Equity Risk Premium (ERP)}: More significantly, the Country Risk Premium (CRP), which is added to the base ERP to calculate the company-specific ERP (Section \ref{sec:erp}), is often directly linked to the sovereign default spread (\(CRP_{\text{Country}} = \text{Default Spread} \times (\frac{\sigma_{Equity}}{\sigma_{Bond}})\)). Therefore, a widening of the sovereign default spread directly increases the CRP and consequently the company-specific ERP. This raises the Cost of Equity (\(C_E\)) for firms with exposure to that country, leading to a higher WACC and lower valuation.
\end{enumerate}
Thus, increased sovereign risk, manifested as wider default spreads, generally translates into higher costs of capital and lower valuations for companies operating within or heavily exposed to those economies.

In summary, the dynamics of both risk-free rates and default spreads are fundamental drivers within the DCF framework. Changes in these rates, driven by macroeconomic policy, economic outlook, or market sentiment, ripple through the cost of equity, cost of debt, and potentially growth assumptions, ultimately impacting the estimated intrinsic value of the firm. The preceding analysis focused primarily on the direct channels through which interest rate changes affect the components of a DCF valuation.\\

In the following chapters, we will broaden the scope to explore the wider macroeconomic consequences that interest rate shifts provoke within the economy, which in turn influence the fundamental assumptions underlying valuation forecasts. We will begin by examining the domestic effects on key economic aggregates such as Aggregated Demand (AD) and its components: consumption, investment, government spending and net exports. Subsequently, our focus will shift to the international dimension, exploring the foreign exchange market as a primary adjustment mechanism and analyzing the resulting implications for international trade dynamics.
